\chapter{Introduction}
\label{ch:intro}

Vehicle registration numbers are the main way to identify a vehicle at a
checkpoint or a traffic junction, and Automatic Number Plate Recognition (ANPR)
systems read them automatically from camera images. Checking plates by hand is
slow and cannot keep up with traffic, so police, toll and parking operators
increasingly rely on cameras that read the plate, look the number up in a
database and raise an alert when a vehicle is stolen, blacklisted or missing
its documents.

Most ANPR systems are designed and tested in clear daylight. Outdoors they meet
haze, fog, rain, glare and low contrast, which blur the plate characters and
break the assumptions on which detectors and character readers were trained. In
our own tests a general-purpose OCR engine read only 16.5--25.2\,\% of plates
exactly in fog and 5.9--11.0\,\% in rain. Kerala makes the problem concrete:
the monsoon brings months of rain, and the hill roads around Ponmudi are in
dense fog for much of the day, yet these are exactly the conditions in which
vehicle checks are most needed. A dependable system therefore has to restore
degraded images, read plates in real time as vehicles pass, and cope with
Indian plate formats, which include one-line and two-line layouts, white,
yellow and red plates, and many state and series codes.

This project develops a \emph{Weather Adaptive ANPR System}. Each camera node
detects plates with YOLOv11, restores weather-degraded frames, and reads the
plate with a recognition network trained on Indian plates. The registration
number is then checked against a central vehicle database and the vehicle is
classified as authorised, non-compliant, unauthorised, blacklisted or stolen.
Alerts appear on a real-time dashboard, and sightings from several camera nodes
are correlated to detect convoy movement and route anomalies. Each node is
designed to run on embedded hardware of the Jetson class.

\section{Problem Statement}

Plate recognition in the field fails in three connected ways. First, the
\emph{image} is degraded: fog lowers contrast and adds a veil of scattered
light, rain adds streaks and spray, and low light adds noise. Second, the
\emph{plate} is small: a plate on a vehicle ten to twenty metres away may be
only 25--40 pixels wide, and two-line plates on two-wheelers and autos are
smaller still. Third, the \emph{reader} is brittle: a network trained on clean
plates misreads the characters it has not seen blurred, repeated digits are
dropped, and letters such as \plate{Q}, \plate{O} and \plate{0} are confused.
The system must handle all three while still running fast enough to show a
plate the moment the vehicle passes.

\section{Objectives}

The main objectives of the project are as follows.
\begin{enumerate}
  \item \textbf{Restoration.} Capture frames and restore plates degraded by
        haze, rain or low contrast so that they stay readable.
  \item \textbf{Detection and recognition.} Detect plates and recognise the
        registration number in real time, with format validation and
        multi-frame confirmation.
  \item \textbf{Verification.} Verify each plate against a central vehicle
        database and classify it as authorised, non-compliant, unauthorised,
        blacklisted or stolen.
  \item \textbf{Alerts and dispatch.} Raise instant alerts on a real-time
        dashboard (vehicle details, image, timestamp) and forward stolen or
        blacklisted vehicles to a dispatch module.
  \item \textbf{Correlation.} Correlate sightings across camera nodes to find
        convoy movement, route anomalies and loitering, and answer
        movement-history queries.
\end{enumerate}
Objectives 1 and 2 are the subject of Phase~1 and of this report. Objectives 3
to 5, together with embedded deployment, make up Phase~2 (Chapter~\ref{ch:phase2}).

\section{Scope of Phase 1}

Phase~1 delivers the vision core of the system, \ie\ everything that turns a
camera frame into a confirmed registration number:
\begin{itemize}
  \item the YOLOv11 plate detector (Section~\ref{sec:yolo});
  \item the weather restoration stage: dehazing with a Mamba-based network,
        its physics-guided extension and a CLAHE route for low contrast, and
        a lightweight deraining network (Section~\ref{sec:restoration});
  \item the LPRNet plate reader with format-constrained decoding
        (Section~\ref{sec:lprnet});
  \item tracking, voting and the real-time pipeline, with tests on unseen
        traffic videos and on foggy video from Ponmudi
        (Sections~\ref{sec:realtime} and~\ref{sec:video}).
\end{itemize}

\section{Contributions}

The main results of Phase~1 are the following.
\begin{itemize}
  \item A plate detector with 99.2\,\% mAP@50 and a plate reader that was
        taken from 0.70 to 0.249 character error rate through synthetic
        pre-training, public data and self-labelled Indian plates, and from
        38 to 41 of 51 exact plates on six unseen videos through two rounds of
        labelled real-video crops.
  \item A real-time pipeline that confirms each plate by voting over frames and
        runs at up to 55 frames per second on 1080p video on a laptop GPU.
  \item A fog-removal network (UVMNet) trained and fine-tuned for plates and
        real fog, and a physics-guided procedure that uses the network only to
        \emph{measure} the fog and then removes it from the real photograph
        by inverting the scattering law.
  \item A controlled study on 60 real foggy photographs showing that, with a
        weather-trained reader, better-looking dehazing does not by itself
        raise the number of plates read, whereas combining frames of the same
        vehicle does.
  \item A lightweight deraining network and an evaluation on synthetic and real
        rain.
\end{itemize}

\section{Organisation of the Report}

Chapter~\ref{ch:lit} reviews the relevant literature. Chapter~\ref{ch:design}
describes the system architecture and the camera node. Chapter~\ref{ch:phase1}
presents the implementation and results of Phase~1. Chapter~\ref{ch:phase2}
lists the work planned for Phase~2, and Chapter~\ref{ch:conclusion} concludes
the report. Appendix~\ref{app:env} lists the datasets and software used.
